\begin{center}
    \large
    
    \begin{singlespace}
        \textbf{\englishTitle{}} \\[0.5cm]
    \end{singlespace}
    
    \begin{singlespace}
        Student : \studentEnName{}  \hspace{1.0cm} 
        % 兩個指導教授要寫 Advisors
        \ifdefined\advisorCnNameB
            Advisors: Dr.\, \advisorEnName \\
            \hspace{6.6cm} Dr.\, \advisorEnNameB  \\
        \else
            Advisor: Dr.\, \advisorEnName \\
        \fi
    \end{singlespace}
    
    \vspace{0.5cm}
    \begin{singlespace}
        \NameofDepartmentInstituteEN\\
        National Yang Ming Chiao Tung University\\[0.2cm]
    \end{singlespace}
    
    \textbf{Abstract} \\[0.5cm]

\end{center}

\normalsize 
  
Knowledge-graph-aware (KG-aware) recommender systems augment graph collaborative filtering with item-side semantics drawn from a knowledge graph (KG); on dense, high-quality KGs they generally improve recommendation accuracy, and the regime is well studied. The setting in which the KG itself is sparse or unreliable---which is the practical default for review-derived KGs, cold-start domains, and privacy-constrained applications---has, by contrast, received little dedicated attention. On a sparse review-derived KG (with on average $2.4$ aspect edges per item after filtering), we observe that four representative KG-aware methods (KGAT, KGCL, MCCLK, and KGRec) are uniformly outperformed by pure LightGCN that does not consult the KG at all. The shared structural cause is that each method injects KG signal into the scoring path as a mandatory component, with no mechanism to disengage or attenuate when the KG is unreliable; KG noise therefore contaminates the gradient flow that produces the collaborative representation.

To address this gap, this thesis proposes RA-GARK (Rationale-Aware Gating over Review-Aspect Knowledge graphs), a dual-view recommendation architecture that treats the KG as an explicitly gated side channel rather than a mandatory pipeline component. Its core design is a local-biased fusion gate: when the user--item interaction graph and the KG are jointly considered, the gate dynamically adjusts the relative weight of the two signals per (user, item) pair, allowing the model to adapt flexibly---amplifying the KG when it is informative for that pair, and falling back to the collaborative signal when it is not, thereby improving the predicted score and the resulting recommendation. The gate's bias is initialized to $+5$, so that training begins as effectively pure LightGCN and the KG channel is opened only when gradients indicate it helps, providing a structural graceful-degradation guarantee. Two further design choices support this gating principle. Each item's KG semantics are stored in $A=4$ \emph{latent} aspect slots (a model hyperparameter, independent of how many raw aspect tags any individual item carries); the slot contents are obtained by a truncated SVD of an IDF-weighted item--aspect co-occurrence matrix, which projects all items into a shared latent semantic subspace and serves as the optimization starting point. A softmax aspect-saliency attention then selects, per (user, item) pair, the most salient latent slot to participate in the global view.

On the sparse review-derived KG benchmark, RA-GARK reaches NDCG@20 = $0.1238$, exceeding the strongest KG-aware baseline (KGRec) by $+13.1\%$ and pure LightGCN by $+5.0\%$. The latter demonstrates that the proposed dynamic weighting mechanism flips the contribution of KG signal from net-negative to net-positive even when the KG itself is unchanged. Component-wise ablations show that replacing softmax with sigmoid attention costs $7.0\%$ NDCG@20, while removing the bias-initialized fusion gate or the KG-SVD initialization each costs $5.3\%$. We additionally report a single-dataset methodological observation that the choice of attention normalization should align with the semantic assumption of the rationale, and explicitly frame this as a hypothesis pending multi-dataset replication rather than a general claim.


\vspace{1cm}

\textbf{Keywords: knowledge-graph-aware recommendation; graceful degradation; rationale-aware learning; fusion gate; sparse knowledge graph; attention normalization; collaborative filtering.} 
